\documentclass[11pt]{article}
% preamble.tex — shared preamble for the Flyxion connective-essay series
% Compile target: XeLaTeX. \input this file, or copy its contents, at the
% top of each essay after \documentclass.

\usepackage{fontspec}
\setmainfont{TeX Gyre Termes} % Times-like serif; falls back gracefully if absent
\usepackage{amsmath,amssymb,amsthm}
\usepackage[margin=1.25in]{geometry}
\usepackage[hidelinks]{hyperref}
\usepackage{microtype}
\usepackage{enumitem}
\usepackage{footmisc}

\newtheorem{claim}{Claim}
\newtheorem{definition}{Definition}
\newtheorem{observation}{Observation}

% --- Corpus vocabulary macros -------------------------------------------
\newcommand{\admissible}{\textsc{admissible}}
\newcommand{\inadmissible}{\textsc{inadmissible}}
\newcommand{\repair}{\textit{repair}}
\newcommand{\continuation}{\textit{continuation}}
\newcommand{\rsvp}{\textsc{rsvp}}

% --- Title block helper ---------------------------------------------------
% Usage: \essaytitle{Title}{Subtitle or tagline}
\newcommand{\essaytitle}[2]{%
  \title{#1\\[0.3em]\large #2}
  \author{Flyxion\\\normalsize Independent Researcher}
  \date{\today}
}

\pagestyle{plain}

\essaytitle{Affinity Maturation as an Admissibility Filter}{Clonal selection, apoptotic repair, and the continuity of immune memory}
\begin{document}
\maketitle
\begin{abstract}
Adaptive immunity is usually narrated as a search problem. In the standard story, clonal diversity supplies variants, antigen selects the fitter ones, and germinal centers iteratively optimize antibody affinity. That account is real and indispensable, but it backgrounds an equally real structure already present in the field's own mechanisms: before affinity can be maximized, lymphocyte variants must first be admissible for further participation. Building on Jerne's 1955 natural-selection theory of antibody formation, Burnet's 1957 clonal selection theory, and the later cellular account of somatic hypermutation and germinal-center selection, this essay argues that affinity is best understood as a ranking internal to a stricter filter. Central and peripheral tolerance establish the initial gate. Dark-zone/light-zone cycling then tests whether mutated B-cell variants remain foldable, non-self-reactive, antigen-competitive, and able to secure T follicular helper help. Most cells die by apoptosis, not because the system failed to search, but because deletion is part of the mechanism that keeps future responses coherent. On this reading, immunological memory is not simply a cache of past winners. It is a structured \continuation\ ledger: the retained set of past-admissible lineages available for re-entry into later responses, sometimes productively and sometimes path-dependently, as in immune imprinting.
\end{abstract}

\section{The anchor: from Jerne and Burnet to the germinal center}
The anchor here is not ``immunity'' in general but a specific lineage of argument and mechanism: Jerne's 1955 proposal that antibodies are not instructed into existence by antigen, Burnet's 1957 clonal selection theory, and the now-standard cellular account of affinity maturation in germinal centers. That lineage matters because it displaced an older instructive picture with a selective one. The field's great achievement was to show that adaptive immunity does not manufacture specificity from scratch after antigen arrives. It begins with a pre-existing repertoire, then differentially expands and preserves the clones that survive selection.

Once somatic hypermutation and germinal-center anatomy were clarified, the story became even more persuasive. B cells proliferate in the dark zone, mutate their immunoglobulin variable-region genes through activation-induced cytidine deaminase, move to the light zone, acquire antigen displayed by follicular dendritic cells, present processed peptide on major histocompatibility complex class II, compete for limited T follicular helper (Tfh) help, and then either die, differentiate, or cycle back for another round. Victora, Nussenzweig, and others helped make the cyclic architecture vivid: the germinal center is not merely a place where affinity increases, but a spatially organized testing regime in which mutation and selection are iterated.

Because this mechanism is so elegant, it is commonly described in optimization language. Clones ``search'' sequence space; antigen and T-cell help ``select'' higher-affinity variants; the response ``improves'' over time. None of that is false. The question is whether it is primary.

\subsection{The unexamined premise}
The unexamined premise is that the decisive explanatory variable is affinity maximization, while the other constraints on B-cell survival are secondary side conditions. Immunology often writes as though the system's central task were to climb an affinity landscape, with deletion, tolerance, and memory treated as auxiliary safeguards or historical residues.

My claim is stronger than the harmless reminder that many constraints matter. It is that the backgrounded constraints are load-bearing. The field's own mechanisms show that affinity only ranks variants \emph{within} a prior admissibility structure. A mutated clone that binds strongly but is catastrophically self-reactive, cannot compete for antigen in the right cellular context, cannot obtain Tfh help, or cannot be stably carried forward into memory is not a failed optimum; it is an \inadmissible\ continuation. The system does not first find the best binder and then later ask whether it can be tolerated and retained. It admits only those variants whose further participation will not break the coherence of future immune behavior.

\subsection{Temporal and constitutive priority}
The priority claim needs to be split in two senses early, otherwise it sounds stronger or stranger than it is.

\paragraph{Temporal priority.}
Some admissibility gates do indeed come earlier in time than affinity maturation. Central tolerance in bone marrow and thymus, peripheral anergy, and receptor editing all operate before or alongside later antigen-driven refinement. In that temporal sense, admissibility often precedes optimization.

\paragraph{Constitutive priority.}
The more important claim is constitutive, not temporal. Even when affinity maturation is already underway, admissibility is what determines what counts as a live candidate for continuation at all. Memory likewise temporally precedes many later recall responses, but the essay's stronger claim is constitutive: the available memory compartment helps determine what future responses can count as admissible, because recall responses are built from lineages that have already survived prior filters. The repertoire is not merely a past cause of later behavior. It helps define the space of continuable options.

\section{Minimum apparatus: admissibility and continuation}
\begin{definition}
For the present essay, a B-cell variant is \admissible\ if it can continue through the next immune checkpoint without importing a contradiction into the response: it must remain functionally expressible as B-cell receptor or antibody, avoid disqualifying self-reactivity, acquire enough antigen under germinal-center conditions to present for Tfh evaluation, and receive sufficient survival signals either to recycle, to differentiate, or to persist as memory. Variants that fail these conditions are \inadmissible.
\end{definition}

Two points follow. First, this is not a re-description of ``high affinity'' in fancier terms. High affinity helps, but admissibility includes properties that affinity alone does not settle. Second, \continuation\ is not an extra module pasted on after selection. It is the immune system's answer to which previously screened lineages remain available for future use. In that limited sense, memory behaves like a ledger: not a perfect archive of the past, and not a record of every competitor, but a retained, actionable set of prior admissions.\footnote{This use of \continuation\ and ``ledger'' is intentionally narrow. It draws only on verified companion essays in this repository: \texttt{theory-rsvp.html}, which reframes physics from ``what moves?'' to ``what is allowed to persist?''; \texttt{theory-holonomy.html}, which treats identity as a loop property and repair as continuation without exact restoration; \texttt{theory-spherepop.html}, which sketches an append-only event-log ontology; and \texttt{theory-unified.html}, which argues that local coherence is easy while global compatibility is hard. Only the idea that persistence is filtered, not automatically given, is borrowed here.}

\begin{claim}
In germinal-center biology, affinity is best understood as a ranking internal to the \admissible\ set, not as the criterion that creates the set.
\end{claim}

\section{Case one: dark-zone/light-zone cycling is a filter before it is an optimizer}
The clearest place to test the claim is the germinal center itself. After activation, B cells enter germinal centers and divide rapidly in the dark zone. Somatic hypermutation then generates daughter cells with altered receptors. Many mutations are neutral or harmful; only some improve binding to the immunizing antigen. The search language fits this stage well, because diversity is being generated. But the decisive events occur when those cells move to the light zone. Quantitatively, this is not a perfectly symmetric shuttle: reviews of germinal-center dynamics describe a rough two-to-one excess of dark-zone over light-zone cells together with substantial cyclic reentry of selected light-zone cells back into the dark zone, so the architecture is a measurable recirculating regime rather than a merely qualitative cartoon.\footnote{Luka Mesin, Jonatan Ersching, and Gabriel D. Victora, ``Germinal center B cell dynamics,'' \emph{Immunity} 45, no. 3 (2016): 471--482.}

There, selection is not performed against an abstract affinity meter. B cells must acquire antigen held on follicular dendritic cells as immune complexes, internalize it, process it, and present derived peptide to Tfh cells. Tfh help is limiting. The selected cell is therefore not merely the one with the best intrinsic equilibrium constant in isolation. It is the one that remains viable in a particular multicellular circuit: receptor function, antigen capture, competition, peptide presentation, and receipt of CD40L- and cytokine-mediated help have to line up.

This is where the admissibility framing adds specificity. A clone can fail in at least four ways that the word ``suboptimal'' blurs together. It may mutate into nonfunctional receptor expression. It may preserve expression but lose enough antigen capture to lose the Tfh competition. It may gain reactivity that risks self-recognition and later pathology. Or it may simply fail to receive survival signals in time. Mechanistically, that last outcome is not just a functional gloss: germinal-center B-cell deletion is well established to run through concrete apoptotic programs, including FAS/FASL-mediated killing and BIM-dependent death when adequate survival help is not secured. In all such cases, apoptosis is not an accidental cost of a search algorithm; it is the local \repair\ operation by which the germinal center removes branches that cannot be coherently carried forward. The system does not restore a pre-mutation state. It instead preserves global tractability by refusing continuation to most candidates.

This becomes especially clear in dark-zone/light-zone recycling. Selected cells often re-enter the dark zone for further proliferation and mutation rather than exiting immediately as plasma cells or memory cells. The point of the cycle is not just incremental affinity gain; it is repeated re-testing under conditions that can disqualify even previously successful lineages. A mutation that raises affinity but destroys breadth, stability, or tolerability is not automatically a winner. Germinal centers are brutal precisely because they are enforcing a moving criterion of future continuability under finite help and real tissue constraints.\footnote{The formal parallel is closer than analogy: Yifan Wu et al., \emph{Safety Does Not Compose: Non-Decaying Loop State for Autonomous LLM Agents}, arXiv:2608.27141 (2026), \url{https://arxiv.org/abs/2608.27141}, prove that a monitor local to a single iteration can become information-theoretically incapable of distinguishing an attack distributed across cycles from benign behavior, and that only a persistent, non-decaying cross-iteration state restores discriminating power. Germinal-center cycling exhibits the structural analogue on the biological side: a mutated variant's fitness cannot be assessed from a single dark-zone/light-zone pass in isolation, because tolerance and antigen-competitiveness are properties of a lineage's history across cycles, not of any one round.}

None of this denies the measured increase in average affinity over time. It explains \emph{how} that increase is permitted to occur without destroying immune usability.

\section{Case two: tolerance mechanisms are even more literal admissibility gates}
The optimization picture looks least complete when one steps back from germinal centers to the wider architecture of adaptive immunity. Burnet's own formulation already treated self-reactive clones as subject to deletion. Modern immunology has filled in the mechanisms in much greater detail.

For B cells in the bone marrow, strong self-reactivity can trigger receptor editing, especially renewed light-chain rearrangement, in an attempt to replace a dangerous receptor with a less self-reactive one. Receptor editing is a \repair\ in the specific sense developed in \textit{Distinction Holonomy}: it is not a restoration of the original receptor, and it does not undo the earlier rearrangement. It substitutes a genuinely different receptor sequence that nonetheless counts as the same cell's continuation, licensed to proceed because the replacement --- not the recovery of what was there before --- satisfies the admissibility criterion going forward. If editing fails, clonal deletion can follow.\footnote{The same graded, non-identical structure appears outside immunology in ordinary tissue repair. Collagen peptides laid down during connective-tissue remodeling do not reconstruct the original fibril network; they form a related but distinct matrix that is judged adequate by mechanical and biochemical criteria, not by literal identity with the tissue that preceded an injury. The comparison is offered only to show that non-identical repair is a general feature of biological continuation, not a special pleading invented for receptor editing.} Other cells enter an anergic state in which they persist but are functionally damped. T cells undergo their own central tolerance regime in the thymus, where cells with dangerous self-reactivity are negatively selected, while peripheral tolerance further restrains escapees. The exact machinery differs between B and T lineages, but the architectural point is shared: a large repertoire is not simply generated and then later optimized for foreign antigen. It is pre-screened so that later responses begin from an already thinned admissible base.

Tolerance therefore shows, in a more literal way than the germinal center does, why admissibility is not optional vocabulary. The system is openly willing to discard receptor configurations \emph{before} any question of high foreign-antigen affinity is answered. Indeed, the classical fear haunting clonal selection theory was autoimmunity: a chronic affliction in which the discarding mechanism itself fails. A receptor with excellent binding properties is not a success if its success destroys the organism whose immune system selected it.

Receptor editing is particularly revealing. It is not simple deletion, and it is not optimization toward the immunizing antigen. It is a constrained attempt to recover an \admissible\ lineage from an \inadmissible\ receptor state. The immune system does not insist on returning to the exact prior sequence; it accepts a non-identical but safer continuation. That is why the language of \repair\ is apt here. The repaired cell is not restored byte-for-byte. It is rendered fit to proceed without smuggling forward the contradiction of self-attack.

\section{Case three: original antigenic sin shows memory as continuation, not merely as best-score retention}
If germinal centers and tolerance show how lineages are admitted or removed, immune imprinting shows why memory cannot be understood as a neutral record of the best current binder. The classical language of original antigenic sin, first made famous in influenza, names the observation that first exposure leaves a durable mark on later responses to related but non-identical antigens. More recent work often prefers the less moralized term ``immune imprinting,'' but the core point is the same: past memory B-cell compartments can dominate recall even when the present antigen has shifted.

That phenomenon is awkward for a simple optimization picture. If the system were only maximizing present binding to the new antigen, one would expect naive clones targeting newly exposed variant epitopes to compete on a relatively clean slate. Yet recall responses are often biased toward previously admitted lineages that recognize conserved or familiar epitopes. Those lineages have practical advantages: they are more numerous, faster to activate, and already licensed by past selection. In influenza, such bias can be protective when conserved epitopes matter, but it can also narrow attention away from antigenic novelty. SARS-CoV-2 work now supports a similarly specific point rather than a vague analogy: Aydillo et al. report that COVID-19 patients' antibody responses bear measurable marks of prior exposure history, so the response to the present infection is not determined solely by the currently encountered antigen.\footnote{Alba Aydillo et al., ``Immunological imprinting of the antibody response in COVID-19 patients,'' \emph{Nature Communications} 12 (2021): 3781.}

This is exactly what it means to say that memory is a \continuation\ ledger. The memory compartment does not merely store the historically best antibody. It stores a portfolio of previously \admissible\ lineages, and the fact of their storage changes the next round's starting conditions. A later response is therefore path-dependent. The past does not merely precede the future; it partly constitutes the repertoire from which the future can be rapidly assembled.

Seen this way, original antigenic sin is not a minor embarrassment to an otherwise perfect optimizer. It is what one should expect from any system whose successful past continuations remain materially available and competitively advantaged in the present. Sometimes that is exactly what host defense needs. Sometimes it is a burden. In either case, memory is acting as selective availability, not as a scalar score.

\section{A disconfirming instance: classic hapten systems and live clonal imaging press back hard}
An honest disconfirming case comes from the very literature that established affinity maturation as a central concept. In classic hapten systems, especially the well-studied NP responses in mice, serum affinity rises over time, particular high-affinity clonotypes are recurrently enriched, and the language of mutation plus selection tracks the data remarkably well. Here the optimization framing is not a straw man. It earns its keep.

Indeed, if one narrows the descriptive window enough---fix the antigen, bracket self-reactivity, ignore repertoire history, and ask only why average affinity increased from one sampling point to the next---then local hill-climbing may be the most economical explanation available. The admissibility reframing does not replace that explanation. It risks emptiness if it merely renames every selection event a ``filter'' without identifying any empirical consequence beyond the already known rise in affinity.

Tas et al. supply a harder challenge still. Using live imaging to track germinal-center clonal dynamics, they found that many clones do not disappear by any clean visible ranking in which lower-affinity variants are simply and legibly removed while higher-affinity variants persist. Instead, much clonal loss appears stochastic, with many lineages going extinct in ways that look closer to drift or chance failure to reenter the cycle than to a crisp admissibility screen.\footnote{Jeroen H. M. Tas et al., ``Visualizing antibody affinity maturation in germinal centers,'' \emph{Science} 351, no. 6277 (2016): 1048--1054.} That matters because it blocks an easy move available to my thesis. If every extinction were read as meaningful exclusion, the admissibility language would be doing too much interpretive work. Some losers may simply be losers by stochastic attrition.

The best reconciliation, if there is one, is weaker than a full vindication. One can say that admissibility sets thresholds within which genuinely stochastic clonal drift still occurs: variants may have to remain foldable, non-self-destructive, antigen-competitive, and help-eligible to stay in play at all, while chance still determines which among the still-viable lineages actually persist. But that is a substantive concession, not a verbal rescue. Tas et al. suggest that germinal centers are not transparently filtering machines. They are mixed regimes in which exclusion and drift can be simultaneously real. If that tension remains partly unresolved, it should remain visible in the essay rather than being explained away.

So where does the reframing add content? It adds content only if it explains why certain apparently promising variants are excluded, why apoptosis is so pervasive, why tolerance architecture must be built into the very possibility of successful selection, and why memory can later bias responses away from the present affinity optimum. In settings where none of those features is at issue, the optimization idiom may simply be the better local description. In settings highlighted by Tas et al., even the filter language itself has to be used with restraint. The right conclusion is therefore not that immunology has been conceptually wrong, but that its optimization story is nested inside a stricter account of what is allowed to survive and be reused, while still leaving room for irreducibly stochastic lineage loss.

\section{When this argument would be wrong or empty}
The central move would be wrong if admissibility collapsed empirically into affinity, such that every criterion invoked here were either reducible to measured antigen binding or epiphenomenal to it. In that case the essay would merely be adding vocabulary to an already sufficient account.

It would also be empty if immunology already treated continuation, tolerance, and deletion as fully co-equal with affinity in its dominant explanatory grammar. But that is not quite the field's rhetoric. Textbooks and reviews routinely state the extra constraints, while still telling the overall story as improvement of binders by selection. The essay depends on that mismatch between acknowledged mechanism and privileged narrative.

Finally, the move survives only if counterexamples are taken seriously. The search for disconfirming cases has therefore extended not only to cases in which affinity maximization seems to do nearly all the explanatory work, as in the classic hapten literature, but also to cases in which even lineage loss is not cleanly interpretable as principled filtering, as in Tas et al.'s live-imaging results. That limits the scope of the claim from both sides. The reframing is strongest when the problem is not merely which clone binds better now, but which lineage can be allowed to remain part of future immunity without undermining tolerance, functional expression, or later response flexibility.

\section{A short speculative coda}
This last suggestion is underdeveloped, but the pattern plausibly recurs because biological systems cannot afford to treat persistence as free. An adaptive immune response is not a one-shot computation whose outputs vanish after evaluation. Its products remain in bodies, tissues, and memory compartments, and those products alter later possibilities. Whenever past winners remain available for reuse, the real problem is not simply optimization over a momentary landscape. It is managing which lineages may continue without making later states incoherent. The germinal center then appears less as a search engine than as a disciplined border crossing between mutation and future admissibility.

\end{document}
